\section*{Questão 1}
\textbf{Enunciado:} Considere uma dada aplicação, na qual a saída do sensor é $-5{,}5\text{ V}$ (contínuo). O controlador que irá receber o sinal do sensor só aceita valores positivos e com amplitude máxima de $3\text{ V}$. Qual estrutura de amplificador pode-se utilizar? Defina valores de resistores e desenhe o sinal de saída.

\subsection*{Resolução Passo a Passo}
\begin{enumerate}
    \item \textbf{Análise do Sinal de Entrada e Desejado na Saída:}
    \begin{itemize}
        \item Entrada: $V_i = -5{,}5\text{ V}$ (sinal contínuo e negativo).
        \item Saída máxima permitida: $V_o \le 3\text{ V}$ (sinal contínuo e positivo).
    \end{itemize}
    
    \item \textbf{Determinação da Estrutura:}
    Como precisamos inverter a polaridade do sinal (de negativo para positivo), devemos utilizar um \textbf{Amplificador Inversor}.

    \item \textbf{Cálculo do Ganho de Tensão ($A_v$):}
    \begin{equation*}
        V_o = A_v \cdot V_i \implies 3\text{ V} = A_v \cdot (-5{,}5\text{ V}) \implies A_v = -\frac{3}{5{,}5} = -\frac{6}{11} \approx -0{,}545
    \end{equation*}

    \item \textbf{Dimensionamento dos Resistores:}
    A equação do ganho para a configuração inversora é:
    \begin{equation*}
        A_v = -\frac{R_f}{R_1} = -\frac{6}{11}
    \end{equation*}
    Escolhendo valores comerciais padronizados para os resistores:
    \begin{itemize}
        \item $R_1 = 11\text{ k}\Omega$
        \item $R_f = 6\text{ k}\Omega$ (ou uma combinação/trimpot de $5{,}6\text{ k}\Omega + 390\ \Omega$)
    \end{itemize}

    \item \textbf{Sinal de Saída:}
    O sinal de saída é uma tensão contínua (DC) estável com valor fixo de $+3{,}0\text{ V}$.
\end{enumerate}



\section*{Questão 2}
\textbf{Enunciado:} Ao se aplicar a tensão $V_i$ (presente na Figura 1 -- com valor de pico de $2{,}5\text{ V}$), em uma dada configuração de amplificador operacional, obteve-se na saída a tensão $V_o$ (Figura 1 com valor de pico de $7{,}5\text{ V}$). Qual a configuração de amplificador operacional que foi utilizada? E qual a relação entre os resistores utilizados?

\subsection*{Resolução Passo a Passo}
\begin{enumerate}
    \item \textbf{Análise da Fase do Sinal (Figura 1):}
    \begin{itemize}
        \item Observa-se que nos instantes em que $V_i$ é positivo, $V_o$ também é positivo.
        \item Os sinais de entrada e saída estão \textbf{em fase} (defasagem de $0^\circ$).
        \item Portanto, a configuração utilizada é o \textbf{Amplificador Não-Inversor}.
    \end{itemize}

    \item \textbf{Cálculo do Ganho de Tensão ($A_v$):}
    \begin{equation*}
        A_v = \frac{V_{o,pico}}{V_{i,pico}} = \frac{7{,}5\text{ V}}{2{,}5\text{ V}} = 3
    \end{equation*}

    \item \textbf{Relação Entre os Resistores:}
    A fórmula do ganho do amplificador não-inversor é:
    \begin{equation*}
        A_v = 1 + \frac{R_f}{R_1} = 3
    \end{equation*}
    \begin{equation*}
        \frac{R_f}{R_1} = 3 - 1 = 2 \implies R_f = 2 \cdot R_1
    \end{equation*}
    \textbf{Conclusão:} O resistor de realimentação ($R_f$) deve possuir o dobro da resistência do resistor ligado à terra ($R_1$).
\end{enumerate}



\section*{Questão 3}
\textbf{Enunciado:} Ao se aplicar a tensão $V_i$ (presente na Figura 2 -- com valor de pico de $2{,}5\text{ V}$), em uma dada configuração de amplificador operacional, obteve-se na saída a tensão $V_o$ (Figura 2 com valor de pico de $5\text{ V}$). Qual a configuração de amplificador operacional que foi utilizada? E qual a relação entre os resistores utilizados?

\subsection*{Resolução Passo a Passo}
\begin{enumerate}
    \item \textbf{Análise da Fase do Sinal (Figura 2):}
    \begin{itemize}
        \item Observa-se que nos instantes em que $V_i$ atinge seu pico positivo ($+2{,}5\text{ V}$), o sinal de saída $V_o$ atinge seu pico negativo ($-5\text{ V}$).
        \item Os sinais estão \textbf{invertidos/fora de fase} (defasagem de $180^\circ$).
        \item Portanto, a configuração utilizada é o \textbf{Amplificador Inversor}.
    \end{itemize}

    \item \textbf{Cálculo do Ganho de Tensão ($A_v$):}
    \begin{equation*}
        A_v = \frac{V_o}{V_i} = -\frac{5\text{ V}}{2{,}5\text{ V}} = -2
    \end{equation*}

    \item \textbf{Relação Entre os Resistores:}
    A fórmula do ganho para a configuração inversora é:
    \begin{equation*}
        A_v = -\frac{R_f}{R_1} = -2
    \end{equation*}
    \begin{equation*}
        \frac{R_f}{R_1} = 2 \implies R_f = 2 \cdot R_1
    \end{equation*}
    \textbf{Conclusão:} O resistor de realimentação ($R_f$) é igual ao dobro do resistor de entrada ($R_1$).
\end{enumerate}

\newpage

\section*{Questão 4}
\textbf{Enunciado:} Um dado sistema de monitoração de dados possui dois sensores distintos. A saída desses dois sensores precisam ser somadas de maneira que nenhum ganho extra seja inserido no sistema. Qual configuração com amplificador operacional pode ser utilizada?

\subsection*{Resolução Passo a Passo}
\begin{enumerate}
    \item \textbf{Identificação do Objetivo:}
    Precisamos realizar a operação matemática de soma entre dois sinais ($v_1$ e $v_2$), de tal forma que a saída reflita a soma com ganho unitário ($|A_v| = 1$), ou seja:
    \begin{equation*}
        v_o = -(v_1 + v_2) \quad \text{ou} \quad v_o = v_1 + v_2
    \end{equation*}

    \item \textbf{Configuração Recomendada:}
    A estrutura ideal é o \textbf{Somador Inversor} (ou Amplificador Somador).
    
    \item \textbf{Dimensionamento dos Resistores:}
    A relação geral da tensão de saída de um somador inversor de duas entradas é dada por:
    \begin{equation*}
        v_o = -\left( \frac{R_f}{R_1} v_1 + \frac{R_f}{R_2} v_2 \right)
    \end{equation*}
    Para garantir que nenhum ganho extra seja introduzido ($|A_v| = 1$), basta escolher todos os resistores iguais:
    \begin{equation*}
        R_1 = R_2 = R_f = R
    \end{equation*}
    Dessa forma:
    \begin{equation*}
        v_o = -(v_1 + v_2)
    \end{equation*}
    \textit{Nota:} Se a inversão de fase não for desejada na aplicação, pode-se colocar um inversor unitário ($R_f/R_1 = 1$) logo em seguida.
\end{enumerate}



\section*{Questão 5}
\textbf{Enunciado:} Considere a configuração com amplificador apresentada na Figura 3. As fontes de entrada $v_1$ e $v_2$ são senoidais, com valores de pico de $3\text{ V}$ e $4\text{ V}$, respectivamente. Desenhe o sinal de saída $v_o$.

\subsection*{Resolução Passo a Passo}
\begin{enumerate}
    \item \textbf{Análise do Circuito da Figura 3:}
    O circuito apresentado é um \textbf{Amplificador Diferencial} (ou Subtrator) básico.
    \begin{itemize}
        \item Resistor do terminal inversor ($R_1$): $5\text{ k}\Omega$
        \item Resistor de realimentação ($R_f$): $5\text{ k}\Omega$
        \item Resistor do terminal não-inversor ($R_3$): $5\text{ k}\Omega$
        \item Resistor de aterramento da entrada não-inversora ($R_4$): $5\text{ k}\Omega$
    \end{itemize}

    \item \textbf{Equação do Amplificador Diferencial:}
    A expressão geral para a saída de um subtrator é:
    \begin{equation*}
        v_o = \left( \frac{R_4}{R_3 + R_4} \right) \left( 1 + \frac{R_f}{R_1} \right) v_2 - \left( \frac{R_f}{R_1} \right) v_1
    \end{equation*}
    Substituindo os valores dos resistores ($R_1 = R_f = R_3 = R_4 = 5\text{ k}\Omega$):
    \begin{equation*}
        v_o = \left( \frac{5k}{5k + 5k} \right) \left( 1 + \frac{5k}{5k} \right) v_2 - \left( \frac{5k}{5k} \right) v_1 = \left( \frac{1}{2} \right) (2) v_2 - (1) v_1 = v_2 - v_1
    \end{equation*}

    \item \textbf{Cálculo da Tensão de Saída $v_o$:}
    Considerando que as fontes $v_1$ e $v_2$ são senoidais em fase:
    \begin{equation*}
        v_1(t) = 3 \sin(\omega t)\text{ V} \quad \text{e} \quad v_2(t) = 4 \sin(\omega t)\text{ V}
    \end{equation*}
    Substituindo na equação do circuito:
    \begin{equation*}
        v_o(t) = 4 \sin(\omega t) - 3 \sin(\omega t) = 1 \sin(\omega t)\text{ V}
    \end{equation*}

    \item \textbf{Características do Sinal de Saída:}
    \begin{itemize}
        \item \textbf{Forma de onda:} Senoidal em fase com as entradas.
        \item \textbf{Amplitude de pico ($V_{o,pico}$):} $1\text{ V}$.
        \item \textbf{Valor pico a pico ($V_{o,pp}$):} $2\text{ V}$ (variando de $-1\text{ V}$ a $+1\text{ V}$).
    \end{itemize}
\end{enumerate}

\newpage

\section*{Questão 6}
\textbf{Enunciado:} Caso ambas as fontes de entrada $v_1$ e $v_2$ da Figura 3, possuíssem um nível CC de $1\text{ V}$, como ficaria a tensão de saída (desenhe)?

\subsection*{Resolução Passo a Passo}
\begin{enumerate}
    \item \textbf{Expressão das Fontes com Nível CC:}
    \begin{equation*}
        v_1(t) = 1 + 3\sin(\omega t)\text{ V} \quad \text{e} \quad v_2(t) = 1 + 4\sin(\omega t)\text{ V}
    \end{equation*}

    \item \textbf{Aplicação na Equação da Figura 3 ($v_o = v_2 - v_1$):}
    \begin{equation*}
        v_o(t) = [1 + 4\sin(\omega t)] - [1 + 3\sin(\omega t)] = (1 - 1) + (4 - 3)\sin(\omega t) = 1\sin(\omega t)\text{ V}
    \end{equation*}

    \item \textbf{Conclusão e Comportamento da Saída:}
    Como o circuito é um amplificador diferencial (subtrator), o nível CC comum de $1\text{ V}$ presente em ambas as entradas é totalmente cancelado (\textit{rejeição em modo comum}). O sinal de saída permanece idêntico ao da Questão 5: uma senóide pura com pico de $1\text{ V}$ e sem deslocamento CC ($V_{CC,sa\text{í}da} = 0\text{ V}$).
\end{enumerate}



\section*{Questão 7}
\textbf{Enunciado:} Considere a Figura 3. Caso a fonte de entrada $v_1$ tivesse nível CC de $1\text{ V}$ e a fonte de entrada $v_2$ de $2\text{ V}$, como ficaria a tensão de saída (desenhe)?

\subsection*{Resolução Passo a Passo}
\begin{enumerate}
    \item \textbf{Expressão das Fontes com Níveis CC Distintos:}
    \begin{equation*}
        v_1(t) = 1 + 3\sin(\omega t)\text{ V} \quad \text{e} \quad v_2(t) = 2 + 4\sin(\omega t)\text{ V}
    \end{equation*}

    \item \textbf{Aplicação na Equação Diferencial ($v_o = v_2 - v_1$):}
    \begin{equation*}
        v_o(t) = [2 + 4\sin(\omega t)] - [1 + 3\sin(\omega t)] = (2 - 1) + (4 - 3)\sin(\omega t) = 1 + 1\sin(\omega t)\text{ V}
    \end{equation*}

    \item \textbf{Sinal de Saída:}
    A saída é uma onda senoidal com amplitude de pico de $1\text{ V}$, deslocada para cima por uma tensão contínua (offset CC) de $+1\text{ V}$. O sinal oscila entre $0\text{ V}$ (mínimo) e $+2\text{ V}$ (máximo).
\end{enumerate}



\section*{Questão 8}
\textbf{Enunciado:} Em quais aplicações deve-se utilizar o amplificador de instrumentação ao invés do amplificador diferencial?

\subsection*{Resolução Passo a Passo}
\begin{itemize}
    \item \textbf{Ponte de Wheatstone e Sensores de Extensometria (Strain Gauges):} Onde as variações de tensão são da ordem de milivolts e a impedância da fonte deve ser preservada sem carregar o sensor.
    \item \textbf{Sinais Biológicos e Médicos (ECG, EEG, EMG):} Em ambientes com altíssimo ruído de modo comum (como o ruído da rede elétrica de 60 Hz) e sinais na faixa de microvolts/milivolts.
    \item \textbf{Aplicações com Fontes de Alta Impedância:} O amplificador de instrumentação possui buffers de entrada não-inversores que oferecem impedância de entrada extremamente alta (na casa de $G\Omega$), diferente do subtrator básico onde a impedância depende dos resistores $R_1$ e $R_3$.
    \item \textbf{Aplicações com Alto Ruído em Modo Comum:} O amplificador de instrumentação oferece um CMRR (Razão de Rejeição em Modo Comum) muito mais elevado e ganho facilmente ajustável por um único resistor ($R_1$).
\end{itemize}

\newpage

\section*{Questão 9}
\textbf{Enunciado:} Considerando o circuito da Figura 4 (Amplificador de Instrumentação de 3 AmpOps), qual é o valor da tensão de saída, se $R_4 = R_3 = R_2 = R = 15\text{ k}\Omega$, $R_1 = 2\text{ k}\Omega$, e a diferença entre as tensões de entrada for de $2\text{ V}$?

\subsection*{Resolução Passo a Passo}
\begin{enumerate}
    \item \textbf{Fórmula Geral do Amplificador de Instrumentação (Figura 4):}
    \begin{equation*}
        v_o = \left( 1 + \frac{2 R_2}{R_1} \right) \cdot \left( \frac{R_4}{R_3} \right) \cdot (v_2 - v_1)
    \end{equation*}

    \item \textbf{Substituição dos Valores:}
    \begin{itemize}
        \item $R_2 = 15\text{ k}\Omega$, $R_1 = 2\text{ k}\Omega$
        \item $R_3 = 15\text{ k}\Omega$, $R_4 = 15\text{ k}\Omega \implies \frac{R_4}{R_3} = 1$
        \item $(v_2 - v_1) = 2\text{ V}$
    \end{itemize}

    \item \textbf{Cálculo do Ganho do Primeiro Estágio ($A_{v1}$):}
    \begin{equation*}
        A_{v1} = 1 + \frac{2 \cdot 15\text{ k}\Omega}{2\text{ k}\Omega} = 1 + \frac{30}{2} = 1 + 15 = 16
    \end{equation*}

    \item \textbf{Cálculo da Tensão de Saída Final ($v_o$):}
    \begin{equation*}
        v_o = 16 \cdot 1 \cdot 2\text{ V} = 32\text{ V}
    \end{equation*}
\end{enumerate}



\section*{Questão 10}
\textbf{Enunciado:} Um sensor fornece um sinal senoidal com amplitude de $10\text{ V}$ de pico ($-10\text{ V}$ a $+10\text{ V}$). Esse sinal deve ser conectado a um conversor A/D cuja faixa de entrada é de $0$ a $5\text{ V}$. Projete um circuito de condicionamento com amplificadores operacionais.

\subsection*{Resolução Passo a Passo}
\begin{enumerate}
    \item \textbf{Mapeamento da Equação de Transferência Linear ($v_o = a \cdot v_i + b$):}
    \begin{itemize}
        \item Para $v_i = -10\text{ V} \implies v_o = 0\text{ V}$
        \item Para $v_i = +10\text{ V} \implies v_o = +5\text{ V}$
    \end{itemize}
    \begin{equation*}
        \text{Ganho } a = \frac{5 - 0}{10 - (-10)} = \frac{5}{20} = 0{,}25 \text{ V/V}
    \end{equation*}
    \begin{equation*}
        \text{Offset } b = 0 - (0{,}25 \cdot (-10)) = +2{,}5\text{ V}
    \end{equation*}
    Logo, a relação desejada é:
    \begin{equation*}
        v_o = 0{,}25 \, v_i + 2{,}5\text{ V}
    \end{equation*}

    \item \textbf{Topologia do Circuito (2 Estágios):}
    \begin{itemize}
        \item \textbf{Estágio 1 (Buffer Não-Inversor):} Garante alta impedância de entrada ($>10\text{ M}\Omega$) para não carregar o sensor.
        \item \textbf{Estágio 2 (Somador Inversor e Inversor com Atenuação e Deslocamento):}
        Usamos um somador inversor: $v_{i2} = -\left( \frac{R_f}{R_{in}} v_i + \frac{R_f}{R_{ref}} V_{ref} \right)$.\\
        Para $V_{ref} = -5\text{ V}$, ajustamos $R_f/R_{in} = 0{,}25$ e $R_f/R_{ref} = 0{,}5$.\\
        Em seguida, um inversor com ganho unitário restaura a fase correta.
    \end{itemize}

    \item \textbf{Dimensionamento dos Componentes:}
    \begin{itemize}
        \item $R_f = 10\text{ k}\Omega$, $R_{in} = 40\text{ k}\Omega$ (atenuador por $0{,}25$)
        \item $R_{ref} = 20\text{ k}\Omega$ com $V_{ref} = -5\text{ V}$ (gera offset de $+2{,}5\text{ V}$)
    \end{itemize}

    \item \textbf{Tensão de Alimentação dos AmpOps:}
    Para suportar os picos de $\pm 10\text{ V}$ sem saturação, alimentamos os AmpOps com $V_{CC} = +15\text{ V}$ e $V_{EE} = -15\text{ V}$.
\end{enumerate}



\section*{Questão 11}
\textbf{Enunciado:} Considere um sensor LDR no lugar de $R_4$ em uma Ponte de Wheatstone com $E_{ex} = 20\text{ V}$, $R_1 = 5\text{ k}\Omega$, $R_2 = 1\text{ k}\Omega$, $R_3 = 1\text{ k}\Omega$.

\subsection*{Resolução Passo a Passo}

\paragraph{(a) Determinação da Curva da Tensão de Saída}
A tensão de saída nos nós centrais da ponte é:
\begin{equation*}
    V_{ponte} = E_{ex} \left( \frac{R_{LDR}}{R_3 + R_{LDR}} - \frac{R_2}{R_1 + R_2} \right) = 20 \left( \frac{R_{LDR}}{1\text{ k} + R_{LDR}} - \frac{1\text{ k}}{6\text{ k}} \right)
\end{equation*}
Lendo os pontos aproximados do gráfico $R_{LDR}$ x Luminosidade:
\begin{itemize}
    \item Para $0{,}4\text{ W/m}^2 \implies R_{LDR} \approx 10\text{ k}\Omega \implies V_{ponte} \approx 20 \left( \frac{10}{11} - \frac{1}{6} \right) = 14{,}85\text{ V}$
    \item Para $0{,}8\text{ W/m}^2 \implies R_{LDR} \approx 6\text{ k}\Omega \implies V_{ponte} \approx 20 \left( \frac{6}{7} - \frac{1}{6} \right) = 13{,}81\text{ V}$
    \item Para $1{,}2\text{ W/m}^2 \implies R_{LDR} \approx 4\text{ k}\Omega \implies V_{ponte} \approx 20 \left( \frac{4}{5} - \frac{1}{6} \right) = 12{,}67\text{ V}$
    \item Para $1{,}6\text{ W/m}^2 \implies R_{LDR} \approx 3\text{ k}\Omega \implies V_{ponte} \approx 20 \left( \frac{3}{4} - \frac{1}{6} \right) = 11{,}67\text{ V}$
\end{itemize}

\paragraph{(b) Condição de Equilíbrio da Ponte}
Para $V_{ponte} = 0$:
\begin{equation*}
    \frac{R_{LDR}}{R_3} = \frac{R_2}{R_1} \implies \frac{R_{LDR}}{1\text{ k}} = \frac{1\text{ k}}{5\text{ k}} \implies R_{LDR} = 200\ \Omega = 0{,}2\text{ k}\Omega
\end{equation*}
A ponte estará em equilíbrio quando a resistência do LDR for $200\ \Omega$ (correspondente a um nível elevado de luminosidade no gráfico).

\paragraph{(c) Amplificador de Instrumentação com Ganho de 2}
Para duplicar a saída da ponte ($A_v = 2$), conectamos um amplificador de instrumentação aos nós da ponte.
\begin{equation*}
    A_v = 1 + \frac{2 R_2}{R_1} = 2 \implies \frac{2 R_2}{R_1} = 1 \implies R_1 = 2 R_2
\end{equation*}
Se fixarmos os resistores de ganho $R_2 = 10\text{ k}\Omega$, então o resistor de ajuste será $R_1 = 20\text{ k}\Omega$.

\paragraph{(d) Tensões de Alimentação}
Como as tensões da ponte chegam a aproximadamente $15\text{ V}$, ao duplicar o sinal, a saída necessita de uma faixa ampla. Deve-se alimentar os amplificadores operacionais com no mínimo $V_{CC} = +32\text{ V}$ (ou fonte simétrica de $\pm 20\text{ V}$) para evitar saturação.

\newpage

\section*{Questão 12}
\textbf{Enunciado:} Explique a diferença entre uma ponte de Wheatstone na configuração simples, meia ponte e ponte completa;

\subsection*{Resolução Passo a Passo}
A principal diferença entre as três configurações reside na quantidade de elementos ativos (sensores resistivos, como \textit{strain gauges} ou LDRs) utilizados nos quatro braços da Ponte de Wheatstone, o que impacta diretamente a sensibilidade do sistema, a linearidade e a compensação de temperatura:

\begin{enumerate}
    \item \textbf{Ponte Simples (Quarto de Ponte / \textit{Quarter-Bridge}):}
    \begin{itemize}
        \item \textbf{Estrutura:} Possui apenas \textbf{1 elemento ativo} (sensor) em um dos braços e 3 resistores fixos de precisão nos demais.
        \item \textbf{Sensibilidade:} É a menor das três configurações. A tensão de saída diferencial é aproximadamente:
        \begin{equation*}
            V_{sa\text{í}da} \approx \frac{E_{ex}}{4} \cdot \frac{\Delta R}{R}
        \end{equation*}
        \item \textbf{Aplicações e Limitações:} Utilizada quando o custo é um fator crítico. Sujeita a erros causados por variações térmicas nos cabos e no próprio sensor, exigindo técnicas adicionais de compensação.
    \end{itemize}

    \item \textbf{Meia Ponte (\textit{Half-Bridge}):}
    \begin{itemize}
        \item \textbf{Estrutura:} Utiliza \textbf{2 elementos ativos} (sensores) e 2 resistores fixos. Tipicamente, um sensor trabalha sob tração ($\Delta R > 0$) e o outro sob compressão ($\Delta R < 0$).
        \item \textbf{Sensibilidade:} Ocorre um ganho de sensibilidade igual ao \textbf{dobro} da ponte simples:
        \begin{equation*}
            V_{sa\text{í}da} \approx \frac{E_{ex}}{2} \cdot \frac{\Delta R}{R}
        \end{equation*}
        \item \textbf{Vantagens:} Apresenta auto-compensação de temperatura nativa (se ambos os sensores forem submetidos às mesmas variações térmicas) e melhora sensivelmente a linearidade do sinal.
    \end{itemize}

    \item \textbf{Ponte Completa (\textit{Full-Bridge}):}
    \begin{itemize}
        \item \textbf{Estrutura:} Todos os \textbf{4 braços} da ponte contêm elementos ativos (sensores em pares opostos de tração e compressão).
        \item \textbf{Sensibilidade:} Oferece a máxima sensibilidade possível, sendo o \textbf{quádruplo} da ponte simples:
        \begin{equation*}
            V_{sa\text{í}da} \approx E_{ex} \cdot \frac{\Delta R}{R}
        \end{equation*}
        \item \textbf{Vantagens:} Excelente rejeição a ruídos térmicos, maior relação sinal-ruído (SNR), máxima linearidade e ausência de efeitos parasitas por cabos longos. É a configuração padrão em células de carga industriais de alta precisão.
    \end{itemize}
\end{enumerate}



\section*{Questão 13}
\textbf{Enunciado:} Quais cuidados deve-se tomar com as fontes de alimentação da ponte de Wheatstone?

\subsection*{Resolução Passo a Passo}
Como a tensão de saída da Ponte de Wheatstone é diretamente proporcional à tensão de excitação ($E_{ex}$), qualquer instabilidade na alimentação reflete-se como erro de medição. Os principais cuidados incluem:

\begin{enumerate}
    \item \textbf{Alta Estabilidade e Regulagem da Tensão de Excitação ($E_{ex}$):}
    A fonte deve ser ultra-estável e apresentar baixo \textit{drift} térmico (variação de tensão com a temperatura), pois qualquer variação percentual em $E_{ex}$ se traduz em um erro percentual equivalente na leitura do sensor.

    \item \textbf{Baixo Nível de Ruído e Ripple:}
    Fontes chaveadas comuns com elevado ruído de alta frequência (\textit{ripple}) podem corromper os pequenos sinais de milivolts da ponte. Deve-se priorizar fontes reguladas lineares de baixo ruído (\textit{Low-Noise / LDO}).

    \item \textbf{Técnica de Sensoriamento Remoto (\textit{4-Wire / 6-Wire Sensing}):}
    Em instalações com cabos longos entre a fonte e a ponte, a resistência dos fios provoca queda de tensão, alterando $E_{ex}$. O uso de linhas adicionais de leitura (\textit{Sense lines}) permite que o regulador monitore e corrija a tensão exatamente nos nós da ponte.

    \item \textbf{Aterramento e Blindagem Adequados:}
    Garantir um ponto de terra limpo (GND de instrumentação) e cabos blindados (par trançado com malha) para evitar que a fonte introduza loops de terra (\textit{ground loops}) ou acoplamento capacitivo/indutivo.

    \item \textbf{Alimentação Pulsada ou CA (em Aplicações Específicas):}
    Para evitar auto-aquecimento dos elementos resistivos da ponte (\textit{self-heating}, que altera a resistência), costuma-se utilizar excitação pulsada ou alternada em vez de contínua permanente.
\end{enumerate}

